% !PRINTED 430
\weberpage{430}

Es bleibt noch übrig, wenn $p$ eine Primzahl und $pp_1$ eine in $n-1$
aufgehende Potenz von $p$ ist, $\varphi(pp_1)$ zu bestimmen.

Ist $\alpha$ eine nicht primitive Congruenzwurzel vom Grade $p^{\pi}$ und
$p^{\pi-1} = p_1$, so ist, wenn $\alpha^h$ die niedrigste nach dem Modul
$n$ mit 1 congruente Potenz von $\alpha$ ist, $h$ ein Theiler von $pp_1$,
d.\,h. eine Potenz von $p$; da $h$ aber kleiner als $pp_1$ sein soll, so
ist es auch ein Theiler von $p_1$, und folglich $\alpha^{p_1} \equiv 1
\pmod{n}$, also: eine nicht primitive Congruenzwurzel des Grades $pp_1$
ist zugleich eine Congruenzwurzel des Grades $p_1$. Umgekehrt ist jede
Congruenzwurzel des Grades $p_1$ zugleich eine imprimitive
Congruenzwurzel des Grades $pp_1$. Da nun die beiden Congruenzen
$x^{pp_1} \equiv 1$, $x^{p_1} \equiv 1$ so viele Wurzeln haben, als ihr
Grad beträgt, so bleiben $pp_1 - p_1$ \gesperrt{primitive Congruenzwurzeln} der
ersten übrig, und es ist also
\weq{13}{\varphi(pp_1) = p_1(p-1).}

Die Function $\varphi(a)$ hat also dieselbe Bedeutung, wie im \S.\,132,
und es ist damit zugleich die Anzahl der primitiven Wurzeln der Primzahl
$n$ gleich $\varphi(n-1)$ gefunden. Wir sprechen also noch den Satz aus:

4.\ \gesperrt{Eine Primzahl $n$ hat $\varphi(n-1)$ primitive Wurzeln.}

Damit ist die Existenz von primitiven Wurzeln für jede Primzahl
nachgewiesen und zugleich ihre Anzahl bestimmt. Zur Auffindung der
primitiven Wurzeln haben wir freilich keine andere allgemeine Methode
als das Probiren, was durch einige Kunstgriffe etwas abgekürzt werden
kann.

Ist $g$ eine primitive Wurzel der Primzahl $n$, so sind die Reste der
Potenzen
\weq{14}{1, g, g^2, g^3 \ldots g^{n-2}}
alle von einander verschieden, da, wenn $g^m \equiv g^{m+\mu}$ wäre,
$g^{\mu} \equiv 1$ sein müsste, was nicht möglich ist, so lange $\mu$
kleiner als $n-1$ ist. Es muss also unter den Resten der Reihe (14),
unter denen der Rest 0 nicht vorkommt, jede der Zahlen
\weq{15}{1, 2, 3 \ldots n-1}
und jede nur einmal enthalten sein, oder mit anderen Worten:

5.\ \gesperrt{Ist $a$ eine durch $n$ nicht theilbare Zahl, so giebt es
eine und nur eine Zahl $\alpha$ aus der Reihe der Zahlen
$0, 1, 2 \ldots n-2$, durch die die Congruenz}
\weq{16}{g^{\alpha} \equiv a \pmod{n}}
%CONT>>
